% ---------------------------------------------------------------------
% Section 6: Conclusion. Drafted by Claude; the outline's "Extensions:
% risk aversion" item is folded in as future work pending material.
% ---------------------------------------------------------------------
\section{Conclusion}\label{sec:conclusion}

Blackwell's comparison of experiments ranks information structures for
all decision problems, but ranks few pairs of information structures.
This paper develops a comparison tailored to monotone POMDPs---the
class of partially observed dynamic programs, common in operations
research, in which value functions and optimal policies are monotone
in likelihood-ratio-ordered beliefs. The LR-better order permits the
garbling that degrades a signal process to depend on the state,
requiring only that the garbling in a lower state likelihood-ratio
dominate the garbling in a higher state. The order is implied by
Blackwell sufficiency, implies the monotone quasi-garbling order that
characterizes the value of information in static monotone decision
problems \citep{kim2023comparing}, and---because likelihood-ratio
dominance survives Bayesian updating and state transitions---delivers
value-function comparisons at every horizon of a monotone POMDP. Our
examples show that the reversely monotone structure of the garbling is
essential: when noise instead intensifies with the state, value
functions can cross and no prior-free ranking exists.

Several extensions merit further study. First, our analysis assumes a
risk-neutral DM maximizing expected discounted rewards; extending the
comparison to risk-averse DMs would require tracking how the LR-better
order interacts with concave utility, paralleling the way
\citet{athey2018value} treat payoff functions with concave incremental
returns in static problems. \Claude{The outline lists ``risk
aversion'' as an extension but no material was provided; this sentence
is a placeholder framing---replace or expand when you decide the
scope.} Second, the exact gap between the LR-better and monotone
quasi-garbling orders in dynamic problems remains to be
characterized: is the stronger order necessary for value comparisons
in monotone POMDPs, or merely sufficient? Third, it would be useful to
develop further families of signal processes---beyond the Normal,
Poisson, and Uniform families presented here---for which the
LR-better comparison can be verified directly from model primitives.
